% Answers to Chapter 11 and to practice test 2, translated from sv/back/facit/11.tex. The chapter's
% review questions are reflective and have no answers.

\facitavsnitt{c11}{Chapter 11}{Review: vectors, functions, trigonometry and logarithms}
\begin{facit}

\svar{1.1} \sv{a} $(-1, 6)$ \sv{b} $(5, 4)$ \sv{c} $(6, 15)$ \sv{d} $(13, 7)$

\svar{1.2} \sv{a} $10$ \sv{b} $13$ \sv{c} $\approx 36.9^{\circ}$ \sv{d} $\approx 247.4^{\circ}$

\svar{1.3} \sv{a} $0$ \sv{b} The vectors are perpendicular, since the scalar product is zero.
\sv{c} $45^{\circ}$

\svar{1.4} \sv{a} $(-0.8, 0.6)$ \sv{b} $(-20, 15)$ \sv{c} $(4, -3)$

\svar{1.5} \sv{a} $(9, -12)\,\text{V}$ \sv{b} $15\,\text{V}$ \sv{c} $\approx -53.1^{\circ}$
\sv{d} The voltages are perpendicular and add geometrically: $\abs{\mathbf{U}}$ is the hypotenuse
of a right-angled triangle, which is shorter than the sum of the other two sides.

\svar{2.1} \sv{a} $D_f = [2, \infty)$, $V_f = [0, \infty)$
\sv{b} $D_f = \mathbb{R} \setminus \{-1\}$, $V_f = \mathbb{R} \setminus \{0\}$
\sv{c} $D_f = \mathbb{R}$, $V_f = (-\infty, 4]$

\svar{2.2} \sv{a} $k = -2\,\Omega/^{\circ}\text{C}$, $m = 200\,\Omega$ \sv{b} $150\,\Omega$
\sv{c} $50^{\circ}\text{C}$ \sv{d} Decreasing, since $k < 0$.

\svar{2.3} \sv{a} $f(x) = 80 \cdot 0.85^x$ \sv{b} $\approx 30$ faults
\sv{c} After $9$ whole months ($x \approx 8.5$).

\svar{2.4} \sv{a} $4\,\text{W}$ \sv{b} $100\,\text{W}$ \sv{c} $V_P = [0, 100]\,\text{W}$
\sv{d} $1.2\,\text{A}$

\svar{2.5} \sv{a} $(x + 3)(x - 3)$ \sv{b} $x = 3$ \sv{c} $x + 3$, where $x \neq 3$
\sv{d} $x = -3$

\svar{3.1} \sv{a} $\frac{3\pi}{4}$ \sv{b} $-\frac{2\pi}{3}$ \sv{c} $150^{\circ}$
\sv{d} $\approx 114.6^{\circ}$

\svar{3.2} \sv{a} $6\,\text{V}$ \sv{b} $\omega = 400\pi\,\text{rad/s}$, $f = 200\,\text{Hz}$
\sv{c} $5\,\text{ms}$ \sv{d} $\approx 0.83\,\text{ms}$; the phase is negative, so the curve is
late.

\svar{3.3} $u(t) = 12\sin\left(1000\pi t + \frac{\pi}{6}\right)\,\text{V}$

\svar{3.4} \sv{a} $8\,\text{V}$ \sv{b} $\approx 5.66\,\text{V}$ \sv{c} $0.64\,\text{W}$

\svar{3.5} \sv{a} $t = 1/1200\,\text{s} \approx 0.83\,\text{ms}$
\sv{b} $t = 5/1200\,\text{s} \approx 4.17\,\text{ms}$ \sv{c} $t = 2.5\,\text{ms}$

\svar{4.1} \sv{a} $7$ \sv{b} $\approx 2.70$ \sv{c} $\approx 0.805$ \sv{d} $101$

\svar{4.2} \sv{a} $2$ \sv{b} $2$ \sv{c} $\log 8$ \sv{d} $3$

\svar{4.3} \sv{a} $40\,\text{dB}$ \sv{b} $28\,\text{dB}$ \sv{c} $\approx 25.1$
\sv{d} $\approx 0.50\,\text{V}$

\svar{4.4} \sv{a} $\approx 14.0\,\text{dB}$ \sv{b} $100\,\text{mW}$ \sv{c} $30\,\text{dBm}$

\svar{4.5} \sv{a} $2.2\,\text{s}$ \sv{b} $\approx 3.07\,\text{V}$
\sv{c} $\approx 1.53\,\text{s}$ \sv{d} $\approx 5.47\,\text{s}$

\end{facit}

\facitavsnitt{od2}{Practice test 2}{Vectors, functions, trigonometry and logarithms}
\begin{facit}

\svar{1} \sv{a} $\mathbf{u}$ is an arrow from the origin to $(1, 2)$ and $\mathbf{v}$ an arrow
from the origin to $(-3, 4)$.
\sv{b} $\abs{\mathbf{u}} = \sqrt{5} \approx 2.24$, $\abs{\mathbf{v}} = 5$
\sv{c} $v_u \approx 63.4^{\circ}$, $v_v \approx 126.9^{\circ}$
\sv{d} $\mathbf{w} = (5, 0)$

\svar{2} $f(x) = x + 1$, where $x \neq 2$.

\svar{3} \sv{a} $\approx 4.14\,\text{V}$
\sv{b} $D_f$: $0 \leq t \leq 20\,\text{s}$; $V_f$: $0.17\,\text{V} \leq u(t) \leq 12\,\text{V}$
\sv{c} $\approx 3.26\,\text{s}$

\svar{4} \sv{a} $u(t) = 5\sin\left(200\pi t + \frac{\pi}{2}\right)\,\text{V}$
\sv{b} A sine curve with the amplitude $5\,\text{V}$ and the period $T = 10\,\text{ms}$, shifted
$2.5\,\text{ms}$ to the left: peak $5\,\text{V}$ at $t = 0$, zero at $2.5\,\text{ms}$, trough
$-5\,\text{V}$ at $5\,\text{ms}$, zero at $7.5\,\text{ms}$ and peak again at $10\,\text{ms}$.

\svar{5} $G_{\text{lin}} = 10^{1.6} \approx 39.8$, that is, about $40$ times.

\end{facit}
